% v6 section. Numbers are generated \vsix macros (scripts/e5_analysis_v6.py via
% scripts/hero_v6.py and scripts/v6_macros.py).
\section{Reviewers and Gates on Violations That Live in State}\label{sec:v6reviewers}

% Generated by scripts/hero_v6.py from promoted v6 evidence; do not edit.
\begin{figure*}[t]\centering
\begin{tikzpicture}[baseline]
\begin{axis}[scale only axis,width=3.75cm,height=3.7cm,tick label style={font=\scriptsize},label style={font=\footnotesize},title style={font=\footnotesize,align=center},xmin=-3,xmax=103,ymin=0.4,ymax=6.0,ytick={5.40,4.40,3.40,2.00,1.00},yticklabels={{Replay, budget spent},{Issuer withdrawn},{Stale revision},{Altered payload},{Injected payload}},xlabel={Unauthorized proposals admitted (\%)},title={\textbf{(a)} Without reasoning, reviewers miss state},yticklabel style={font=\tiny,align=right},xmajorgrids,axis y line*=left,clip=false,legend style={font=\tiny,at={(0.5,-0.34)},anchor=north,draw=none,fill=none,legend columns=1,row sep=-1.5pt},legend cell align=left]
\addplot[only marks,mark=triangle*,mark size=2pt,draw=endored,fill=endored,line width=0.6pt] coordinates {(32.43,5.650) (25.53,4.650) (24.49,3.650) (0,2.250) (5,1.250)};\addlegendentry{same-model reviewer, rules}
\addplot[only marks,mark=diamond,mark size=2pt,draw=warnamber,fill=white,line width=0.6pt] coordinates {(0,5.480) (0,4.480) (0,3.480) (0,2.080) (5,1.080)};\addlegendentry{same-model reviewer, rules + reasoning}
\addplot[only marks,mark=*,mark size=2pt,draw=exoblue,fill=exoblue,line width=0.6pt] coordinates {(0,5.320) (0,4.320) (0,3.320) (0,1.920) (0,0.920)};\addlegendentry{gate reading sink state}
\addplot[only marks,mark=o,mark size=1.4pt,draw=structgray!80] coordinates {(32.43,5.150) (25.53,4.150) (24.49,3.150) (33.33,1.750) (60,0.750)};\addlegendentry{gate with that clause unwritten (by construction)}
\draw[structgray!60,densely dashed] (axis cs:-3,2.70) -- (axis cs:103,2.70);
\node[font=\tiny\itshape,text=slate,anchor=south east] at (axis cs:103,5.65) {violation lives in state};
\node[font=\tiny\itshape,text=slate,anchor=north east] at (axis cs:103,2.65) {visible in the proposal};
\end{axis}\end{tikzpicture}
\hspace{1.4cm}
\begin{tikzpicture}[baseline]
\begin{axis}[scale only axis,width=3.75cm,height=3.7cm,tick label style={font=\scriptsize},label style={font=\footnotesize},title style={font=\footnotesize,align=center},xmin=-3,xmax=103,ymin=-3,ymax=103,grid=major,clip=false,xlabel={Benign tasks completed (\%)},ylabel={Attacks that succeed (\%)},title={\textbf{(b)} A specification stops what it names},legend style={font=\tiny,at={(0.5,-0.34)},anchor=north,draw=none,fill=none,legend columns=2,row sep=-1.5pt,column sep=3pt},legend cell align=left]
\addplot[only marks,mark=square*,mark size=2.3pt,draw=structgray,fill=structgray,error bars/.cd,x dir=both,x explicit,y dir=both,y explicit,error bar style={structgray,line width=0.45pt}] coordinates {(43.24,56.76) += (17.27,11.48) -= (16.15,12.04)};\addlegendentry{Unprotected}
\draw[structgray,line width=0.6pt] (axis cs:43.24,56.76) -- (axis cs:43.24,83.78);
\addplot[forget plot,only marks,mark=square,mark size=2.3pt,draw=structgray,fill=white,error bars/.cd,y dir=both,y explicit,error bar style={structgray,line width=0.45pt}] coordinates {(43.24,83.78) += (0,7.546) -= (0,10.4)};
\addplot[only marks,mark=diamond*,mark size=2.3pt,draw=warnamber!85!black,fill=warnamber!85!black,error bars/.cd,x dir=both,x explicit,y dir=both,y explicit,error bar style={warnamber!85!black,line width=0.45pt}] coordinates {(32.43,2.703) += (17.35,6.722) -= (14.42,2.374)};\addlegendentry{Progent (per call)}
\draw[warnamber!85!black,line width=0.6pt] (axis cs:32.43,2.703) -- (axis cs:32.43,4.054);
\addplot[forget plot,only marks,mark=diamond,mark size=2.3pt,draw=warnamber!85!black,fill=white,error bars/.cd,y dir=both,y explicit,error bar style={warnamber!85!black,line width=0.45pt}] coordinates {(32.43,4.054) += (0,7.34) -= (0,3.21)};
\addplot[only marks,mark=*,mark size=2.3pt,draw=exoblue,fill=exoblue,error bars/.cd,x dir=both,x explicit,y dir=both,y explicit,error bar style={exoblue,line width=0.45pt}] coordinates {(24.32,0) += (16.87,4.863) -= (12.55,0)};\addlegendentry{Strict \(S\)}
\draw[exoblue,line width=0.6pt] (axis cs:24.32,0) -- (axis cs:24.32,71.62);
\addplot[forget plot,only marks,mark=o,mark size=2.3pt,draw=exoblue,fill=white,error bars/.cd,y dir=both,y explicit,error bar style={exoblue,line width=0.45pt}] coordinates {(24.32,71.62) += (0,9.877) -= (0,11.67)};
\addplot[only marks,mark=pentagon*,mark size=2.3pt,draw=exoblue,fill=exoblue,error bars/.cd,x dir=both,x explicit,y dir=both,y explicit,error bar style={exoblue,line width=0.45pt}] coordinates {(32.43,10.81) += (17.35,9.384) -= (14.42,6.026)};\addlegendentry{Delegated}
\draw[exoblue,line width=0.6pt] (axis cs:32.43,10.81) -- (axis cs:32.43,70.27);
\addplot[forget plot,only marks,mark=pentagon,mark size=2.3pt,draw=exoblue,fill=white,error bars/.cd,y dir=both,y explicit,error bar style={exoblue,line width=0.45pt}] coordinates {(32.43,70.27) += (0,10.07) -= (0,11.75)};
\addplot[only marks,mark=triangle*,mark size=2.3pt,draw=exoblue!60!black,fill=exoblue!60!black,error bars/.cd,x dir=both,x explicit,y dir=both,y explicit,error bar style={exoblue!60!black,line width=0.45pt}] coordinates {(40.54,0) += (17.36,4.863) -= (15.79,0)};\addlegendentry{Oracle (solution allowlist)}
\draw[exoblue!60!black,line width=0.6pt] (axis cs:40.54,0) -- (axis cs:40.54,0);
\addplot[forget plot,only marks,mark=triangle,mark size=2.3pt,draw=exoblue!60!black,fill=white,error bars/.cd,y dir=both,y explicit,error bar style={exoblue!60!black,line width=0.45pt}] coordinates {(40.54,0) += (0,4.863) -= (0,0)};
\node[font=\tiny,text=slate,anchor=north east,align=right] at (axis cs:103,103) {filled: AgentDojo's injections\\open: injections inside S};
\end{axis}\end{tikzpicture}
\caption{\textbf{A same-model reviewer misses violations in state unless it reasons at length; a written specification stops what it names, at a cost.} \textbf{(a)} Qwen3-8B red-team proposals under 8 predeclared mutation strategies, each distinct proposal counted once (no interval: proposals within a case are not independent). The reviewer is the same checkpoint, shown the rules and the state the gate reads. Without reasoning it admits 36/133 of unauthorized proposals whose violation lives in sink state, more often than 2/76 of those visible in the proposal (by case, 36/36 against 2/24). With reasoning (median 637 generated tokens per review, against 6) it admits 0/133 and 2/76. The gate admits 0; it enforces the written predicate, and its Rust kernel agreed with the Python reference on all 217 valid proposals of both sample sets (0 disagreements). Open circles are by construction: the share of each family's violations that break only that family's clause, which a gate without that clause would admit. \textbf{(b)} AgentDojo banking and Slack (37 user tasks), Qwen3-8B, AgentDojo's graders, every episode counted (95\% Clopper--Pearson; the two attacked episodes of a task are not independent, so attack intervals are too narrow). A provenance specification written without the reference solutions (strict \(S\)) completes 9/37 benign tasks; unprotected, 16/37; the oracle, an allowlist copied from each task's solution and so a reference rather than a deployable specification, 15/37. AgentDojo's injections succeed in 0/74 strict and 42/74 unprotected episodes; injections that name only principals already in the user's state (written by us) succeed in 53/74 strict, 52/74 delegated, 62/74 unprotected, 3/74 Progent and 0/74 oracle episodes. Local quantized open-weight models on CPU at temperature 0; no frontier model.}\label{fig:v6spec}
\end{figure*}


A common alternative to mediation is review: a second model, often the same one, reads each proposed action and
approves or denies it~\cite{control}. Review sees the proposal. Many unauthorized actions, however, are
unauthorized only because of the state they would act on: a payment that was legitimate once and is now a replay, a
grant whose issuer has since been withdrawn, an edit against a revision that has moved. Such a proposal can be
byte-identical to one the policy approved. We test whether a reviewer that is given the rules, the visible state and
time to reason catches these, and what a gate that reads state does instead.

\para{Setup (E5).} Sixty cases in five families (twelve each) each pose a pre-state, a policy with six written
clauses (exact match to an approved proposal, cumulative budget, live issuers, current revision, expiry, scope) and an
objective. Three families are \emph{state-dependent}: their violations break a clause that can only be checked against
sink state (replay after the budget is spent, issuer withdrawn, stale revision). Two are \emph{proposal-visible}: an
altered payload, and a payload altered at the request of an injected note that the proposer, not the reviewer, reads.
The same checkpoint proposes candidates under eight predeclared mutation strategies and reviews them in three arms:
the prior work's minimal prompt (R0), the rules stated (R1), and the rules with reasoning enabled (R2, 4096 tokens; a
truncated review is a deny). Every valid candidate also goes through the real broker and sink at the case's pre-state
(G). Plan, prompts, clauses, parse rules and hypotheses were committed before data
(three pre-data amendments). We count each distinct proposal once and
also report the share of cases with any breach, because repeated samples of a deterministic review are not independent.

\para{Result.} The first predeclared hypothesis (H\textsubscript{E5a}) held on both checkpoints: without reasoning
(R0 and R1), reviewers admitted violations that live in state more often than violations visible in the proposal. With the rules stated (R1), the
\vsixEFiveCkpt{} reviewer admitted \vsixReviewerState{} unauthorized proposals whose violation lives in state and
\vsixReviewerVisible{} whose violation is visible (Figure~\ref{fig:v6spec}a), and it admitted at least one
state-dependent violation in \vsixReviewerStateCases{} cases against \vsixReviewerVisibleCases{} proposal-visible
cases. In the state-dependent families the violating proposal is often an exact copy of an approved one, and a
reviewer that answers in a few tokens approves it. Reasoning changed this. The prediction that reasoning would reduce
but not eliminate state-dependent misses (H\textsubscript{E5b}) held on Qwen3-4B and failed on \vsixEFiveCkpt: with
reasoning enabled (R2), the \vsixEFiveCkpt{} reviewer admitted \vsixReasonerState{} state-dependent and
\vsixReasonerVisible{} visible violations, at a median of \vsixReasonerTokens{} generated tokens per review instead of
\vsixReviewerTokens. Every arm on both checkpoints is in Table~\ref{tab:v6e5} (Appendix~\ref{app:v6e5}). The gate
admitted \vsixGateBreaches{}. That zero is the gate enforcing its written predicate, not a finding: on all
\vsixGateUnique{} distinct valid proposals its Rust implementation agreed with an independent Python statement of the
same clauses (\vsixMirrorDisagreements{} disagreements). \vsixGateReplaySentence

\para{What this shows.} A reviewer can check state only if it is shown the state, and E5 shows it everything the gate
reads: one moment, at most one prior record. Under those conditions a same-model reviewer that reasons at length did
as well as the gate on these cases, and one that does not reason missed a quarter or more of the violations that live
in state. Three differences remain. Two are measured here. The gate's result does not depend on the checkpoint or
the prompt, while the same model's reviewer went from \vsixReviewerState{} to \vsixReasonerState{} on the same
proposals when reasoning was switched on; and the gate's check costs a database transaction where the reviewer's
costs hundreds of generated tokens. The third is not measured in E5: the gate checks at commit, while a reviewer
checks a snapshot before it, so state that changes in between, such as a budget spent by a concurrent call
(Appendix~\ref{sec:systems}), is invisible to the reviewer.

\para{What a gate misses instead.} Remove one clause and the gate admits exactly the violations that break only that
clause. This holds by construction, and Figure~\ref{fig:v6spec}a draws it as such. What the experiment measures is how
much each family rests on a single clause. \vsixSingleClauseSentence{} This is the paper's central hypothesis, which E5
illustrates and does not test: mediation does not make safety independent of the proposer. It moves the failure from
a reviewer's judgment, which here varied with prompt and budget, to a written specification, whose gaps can be listed,
tested and argued about.

\para{Limits.} The reviewers are local open-weight checkpoints at temperature 0 reviewing their own proposals; a
stronger or differently trained monitor may do better, and control protocols that resample or
audit~\cite{bhatt2025ctrlz} are a different defence. The proposer is told the adversarial objective in plain text. The
families differ in more than where the violation lives, so the comparison shows that these state-dependent cases are
hard for a reviewer without reasoning, not that state-dependence alone makes them hard. Only eight distinct proposals
were authorized, so the data cannot estimate how often a reviewer denies legitimate work, and a reviewer that denied
everything would also score zero. \vsixRtwoTruncated{} R2 reviews reached the token cap and count as denials, as the
plan specified; this can only lower R2's admissions.
